%==============================================================
% appendix_foundations.tex —— 附录 A：基础条目速查
%==============================================================
% 定位（2026-11-04）: **不算篇目**。它不是 A ≅ B，故不进 README §2 的六节点，
%   也不触发 B4。它是 BR0–BR4 那两行「这里默认你已经会」的展开。
% 准入规则（与正文的闭环）:
%   · 只收 BR0–BR4 的前置行里点过名的条目；前置没点名的，不进本附录；
%   · 前置点名过的，本附录必须有 —— 缺一条就是闭环破了。
% 格式规则: 每条固定六字段（作用 / 解决什么 / 看到就想到 / 必须会 / 最容易混 / 自检），
%   每条 ≤ 9 行。超过说明在写教材，删。
% 纪律: 不推导、不举例超过一句、不给其它册子的指针。只回答"它是干什么用的、什么时候该想起它、
%   怎么在 3 分钟内知道自己会不会"。**自检不给答案**（2026-11-04）：给了答案，
%   "我会"就退化成"我看得懂答案"，自评立刻失效。
%==============================================================

\appendix
\section{基础条目速查：看到什么，想到什么}
\label{sec:appendix}

\noindent\small
本附录**不是篇目**，是 BR0--BR4 开头那两行「这里默认你已经会」的展开。
用法只有一种：正文某一句卡住时，回到这里查「它默认我会什么」。

每条回答四件事——它是干什么用的、在解决什么问题、**出现什么信号就该想起它**、哪些必须会；
外加一行最常见的混法，和**一条 3 分钟自检**。
自检的判据是过程而不是结果：**3 分钟内能写出完整步骤 = 达标；卡在某一步 = 未达标，
卡住的那一步就是你要补的那块**。自检不给答案——给了答案，自评就变成「我看得懂答案」。\normalsize

\subsection{线性：把对象换成数组}
\label{app:lin}

\subsubsection*{线性空间与基}
\textbf{作用}：让「对象」能写成一列数，从而能用矩阵算。\\
\textbf{解决什么}：同一件事换一组基，数组变了、事没变——所有不变量（谱、迹、行列式）都从这来。\\
\textbf{看到就想到}：「坐标」「分量」「换系」「表示」；同一个对象有两种写法，而某种写法下运算特别简单。\\
\textbf{必须会}：基与维数、坐标、基变换、子空间、秩与解空间维数。\\
\textbf{最容易混}：换基不是换对象。换的是数组，不是那个东西。\\
\textbf{自检}：把 $1+x+x^{2}$ 在基 $\{1,x,x^{2}\}$ 与基 $\{1,\,x-1,\,(x-1)^{2}\}$ 下的坐标都写出来，
并写出两组基之间的变换矩阵。

\subsubsection*{内积、范数、正交}
\textbf{作用}：给出「长度」和「角度」，于是才有「最近」「垂直」「投影」。\\
\textbf{解决什么}：在一堆候选里挑最近的一个；把误差压成一个数；判断两个方向是否无关。\\
\textbf{看到就想到}：「最像」「误差最小」「夹角/相似度」「正交」「能量」。\\
\textbf{必须会}：内积与范数、Cauchy--Schwarz、Gram 矩阵与夹角、正交归一化。\\
\textbf{最容易混}：$L^{2}$ 内积是积分不是求和；「正交」永远相对于某个内积而言。\\
\textbf{自检}：算 $\langle(1,2,3),(0,1,1)\rangle$ 与两者的夹角；再把 $(1,1)$ 对 $(1,0)$ 正交归一化。

\subsubsection*{特征值与对角化}
\textbf{作用}：找一组让算子「不耦合」的方向。\\
\textbf{解决什么}：在 $Ax=\lambda x$ 的方向上矩阵只是缩放；对角化之后 $A^{k}$、$e^{At}$、稳定性全变成逐分量的事。\\
\textbf{看到就想到}：「长期行为」「哪个方向被放大/压缩」「解耦」「主模态」「共振」。\\
\textbf{必须会}：特征方程、可对角化的条件、实对称矩阵可正交对角化、Jordan 型（认得即可）。\\
\textbf{最容易混}：可对角化不是普遍情形；实矩阵的特征值可以是复数（那就是旋转加缩放）。\\
\textbf{自检}：求 $\begin{pmatrix}2&1\\1&2\end{pmatrix}$ 的特征值与特征向量，写出 $P$ 使 $P^{-1}AP$ 为对角。

\subsubsection*{正交投影}
\textbf{作用}：把对象压进「允许的集合」，并且改动最小。\\
\textbf{解决什么}：约束下的最优解；最小二乘；扔掉不该有的分量。\\
\textbf{看到就想到}：「约束」「允许集合」「残差」「拟合」「只保留某一类成分」。\\
\textbf{必须会}：$P=A(A^{H}A)^{-1}A^{H}$（投影到列空间）、$I-P$ 是残差投影、判据 $\langle x-Px,c\rangle=0$。\\
\textbf{最容易混}：投影 $\ne$ 测地线；允许集合不是子空间时，「最近点」不是一次线性运算。\\
\textbf{自检}：把 $b=(1,2,2)$ 投影到 $\mathrm{span}\{(1,0,0),(0,1,0)\}$，再验证残差与这两个向量都正交。

\subsection{分析：把变化率与误差讲清楚}
\label{app:ana}

\subsubsection*{多元微分：偏导、Jacobian、Hessian、链式法则}
\textbf{作用}：把一元的「变化率」搬到多个变量上。\\
\textbf{解决什么}：多变量函数的线性化；复合函数的导数；极值的二阶判据。\\
\textbf{看到就想到}：「敏感度」「线性化」「误差传播」「切平面」「最优化」。\\
\textbf{必须会}：偏导、梯度、Jacobian（一阶，矩阵）、Hessian（二阶，对称矩阵）、$J_{f\circ g}=J_fJ_g$、Hessian 定号判极值。\\
\textbf{最容易混}：Jacobian 是矩阵不是数；Hessian 的对称性要二阶偏导连续才成立。\\
\textbf{自检}：$f(x,y)=(x^{2}y,\ \sin(xy))$，写出 $J_f$ 并求它在 $(1,1)$ 的值；
再用 $J_{f\circ g}$ 写一遍复合函数的链式法则。

\subsubsection*{泰勒展开}
\textbf{作用}：把函数在一点换成一串多项式，于是「局部 $=$ 线性 $+$ 二阶修正 $+\cdots$」。\\
\textbf{解决什么}：估误差、定收敛阶、证极限、把非线性问题局部线性化。\\
\textbf{看到就想到}：「小量」「近似」「一阶项 / 二阶项」「误差是 $O(h^{2})$」。\\
\textbf{必须会}：一阶与二阶展开；$e^{x}$、$\sin x$、$(1+x)^{\alpha}$、$\log(1+x)$；余项是什么意思。\\
\textbf{最容易混}：精度看余项，不看写了几项。\\
\textbf{自检}：写 $e^{x}$、$\sin x$、$\ln(1+x)$ 到三阶；用二阶展开估 $\sqrt{1.02}$，并说出误差是几阶。

\subsubsection*{分部积分与散度定理}
\textbf{作用}：把导数从一个因子搬到另一个因子上。\\
\textbf{解决什么}：降阶（二阶导降成一阶）；把边界条件搬出来；推出弱形式与能量估计。\\
\textbf{看到就想到}：「二阶导不好处理」「边界项」「弱解」「能量法」。\\
\textbf{必须会}：$\int u'v=-\int uv'+\text{边界}$；多维版即散度定理 / Green 公式；带度量与带权的版本认得即可。\\
\textbf{最容易混}：漏掉边界项——大部分「证明写错了」都错在这里。\\
\textbf{自检}：算 $\int_0^1 xe^{-x}dx$；再把 $\int_\Omega u\,\Delta v$ 分部一次，把边界项单独写出来。

\subsubsection*{变分：驻值与约束}
\textbf{作用}：把「解方程」换成「找一个函数，使某个量取驻值」。\\
\textbf{解决什么}：最短路径、最小能量、平衡态；也是有限元与 Euler--Lagrange 的来源。\\
\textbf{看到就想到}：「最短 / 最省 / 最小作用量」「能量最小」「在一族曲线里挑一条」。\\
\textbf{必须会}：$L[\gamma]=\int\lVert\dot\gamma\rVert$ 的驻值条件；带约束（如 $\lVert\gamma\rVert=1$）时乘子进方程；Euler--Lagrange 的一句结论。\\
\textbf{最容易混}：驻值 $\ne$ 极小（可能是极大或鞍点）；变分是对「函数」求导，不是对 $t$。\\
\textbf{自检}：对 $L[\gamma]=\int_0^1\lVert\dot\gamma\rVert\,dt$ 在 $\lVert\gamma\rVert=1$ 下写出驻值条件，
并指出乘子是从哪一步进来的。

\subsection{方程：从变化率到将来}
\label{app:ode}

\subsubsection*{ODE（常微分方程）}
\textbf{作用}：把「变化率已知」翻成「将来会怎样」。\\
\textbf{解决什么}：给定 $\dot x=f(t,x)$ 与起点，问 $x(t)$ 是多少；以及长期行为——收敛、发散、还是周期。\\
\textbf{看到就想到}：出现「变化率 $=$ 当前状态的函数」；「速度 / 增长 / 衰减 / 追及」；
某个量求导之后仍留在方程里；PDE 沿一条曲线限制下来。\\
\textbf{必须会}：可分离变量、一阶线性（积分因子）、二阶常系数线性（特征方程）、
$\ddot x=-x$ 与 $\ddot x=-\omega^{2}x$、初值问题的解；
「解的存在唯一靠 Lipschitz」这句结论认得即可。\\
\textbf{最容易混}：以为 ODE 都该有闭式解——绝大多数没有，实际要的是相图、稳定性与数值解。\\
\textbf{自检}：各解一个——$\dot x=2x$；$\dot x+x=\sin t$；$\ddot x+4x=0$ 带 $x(0)=1,\ \dot x(0)=0$。

\subsubsection*{简谐振子}
\textbf{作用}：所有「被拉回平衡」的线性系统的原型。\\
\textbf{解决什么}：$\ddot x+\omega^{2}x=0$ 的解就是 $\cos,\sin$；阻尼、受迫、共振都是它的变形。\\
\textbf{看到就想到}：回复力、弹簧、小角度摆、任何「偏离一点就被拽回来」；
以及**任何方程在平衡点附近的线性化**。\\
\textbf{必须会}：$\ddot x=-x$、$\ddot x=-\omega^{2}x$ 的解；特征方程法；欠阻尼 / 临界 / 过阻尼三种情形。\\
\textbf{最容易混}：解不一定是正弦——指数与双曲是同一方程另外的参数区。\\
\textbf{自检}：写 $\ddot x+\omega^{2}x=0$ 的通解；再对 $\ddot x+2\zeta\omega\dot x+\omega^{2}x=0$
说清 $\zeta<1$、$\zeta=1$、$\zeta>1$ 三种情形的解长什么样（不必定系数）。

\subsubsection*{矩阵指数 $e^{At}$}
\textbf{作用}：线性 ODE 的解算子。\\
\textbf{解决什么}：$\dot x=Ax$ 从 $x_{0}$ 出发走时间 $t$；也是离散迭代 $x_{k+1}=Ax_{k}$ 的连续版。\\
\textbf{看到就想到}：常系数线性系统、状态转移、旋转的叠加（$A$ 反对称时 $e^{At}$ 就是旋转）。\\
\textbf{必须会}：级数定义；可对角化时 $e^{At}=Pe^{\Lambda t}P^{-1}$；$A$ 反对称时与 Rodrigues 公式同源。\\
\textbf{最容易混}：$e^{A+B}\ne e^{A}e^{B}$，除非 $[A,B]=0$。\\
\textbf{自检}：$A=\begin{pmatrix}0&-1\\1&0\end{pmatrix}$，写出 $e^{At}$ 并说明它是旋转；
再举一组 $2\times2$ 的 $A,B$ 使 $e^{A+B}\ne e^{A}e^{B}$。

\subsection{变换：换一组基让算子变简单}
\label{app:trans}

\subsubsection*{复指数与欧拉公式}
\textbf{作用}：把「旋转」写成「乘一个数」，于是三角函数变成指数。\\
\textbf{解决什么}：旋转叠加 $=$ 指数相加；微分方程的特征根；傅里叶变换的基。\\
\textbf{看到就想到}：「相位」「旋转」「周期」、$\cos$ 与 $\sin$ 总是成对出现、「相量」。\\
\textbf{必须会}：$e^{i\theta}=\cos\theta+i\sin\theta$；$\lvert e^{i\theta}\rvert=1$；乘法对应角度相加。\\
\textbf{最容易混}：$e^{i\theta}$ 就是单位圆上的一个点，不是「虚数」那种玄学。\\
\textbf{自检}：写出 $e^{i\pi/3}$ 的实部与虚部；算 $e^{i\pi/3}\cdot e^{i\pi/6}$ 并说明角度相加；
把 $\cos3\theta$ 用 $\cos\theta$ 表示。

\subsubsection*{傅里叶级数与傅里叶变换}
\textbf{作用}：换一组基，让「微分」变成「乘 $i\omega$」。\\
\textbf{解决什么}：微分方程变代数方程；看信号由哪些频率组成；卷积变乘法。\\
\textbf{看到就想到}：「周期」「频谱」「滤波」「解 PDE 的常用手法」「在两个域之间来回」。\\
\textbf{必须会}：级数系数公式、变换与逆变换、$\widehat{f'}=i\omega\hat f$、Parseval；采样定理认得即可。\\
\textbf{最容易混}：级数（周期、离散谱）与变换（非周期、连续谱）是同一件事的两个版本。\\
\textbf{自检}：写出周期 $2\pi$ 的方波（取 $\pm1$）的傅里叶系数；
再用 $\widehat{f'}=i\omega\hat f$ 说明「微分方程为什么在频域变成代数方程」。

\subsubsection*{卷积}
\textbf{作用}：「先各自作用、再叠加」的统一写法。\\
\textbf{解决什么}：线性系统的响应；独立随机变量之和的分布；滤波。\\
\textbf{看到就想到}：「滑动平均」「冲激响应」「独立随机变量相加」「乘积出现在另一个域」。\\
\textbf{必须会}：$(f*g)(t)=\int f(\tau)g(t-\tau)\,d\tau$；卷积定理 $\widehat{f*g}=\hat f\,\hat g$。\\
\textbf{最容易混}：卷积不是逐点相乘；$f*g$ 通常比 $f$、$g$ 都光滑。\\
\textbf{自检}：算 $\mathbf{1}_{[0,1]}*\mathbf{1}_{[0,1]}$（应当得到一个三角形），并指出它在哪个域变成乘积。

\subsection{概率：把分布压成几个数}
\label{app:prob}

\subsubsection*{期望、方差、独立}
\textbf{作用}：把一个分布压成几个数。\\
\textbf{解决什么}：平均行为、波动大小、两件事是否互相影响。\\
\textbf{看到就想到}：「平均」「波动」「互不影响」「重复试验」。\\
\textbf{必须会}：$\mathbb{E}$、$\mathrm{Var}(X)=\mathbb{E}X^{2}-(\mathbb{E}X)^{2}$、$\mathrm{Var}(X+Y)$、独立 $\Rightarrow$ 期望可乘、协方差。\\
\textbf{最容易混}：不相关 $\ne$ 独立（高斯情形除外）；$\mathbb{E}[XY]=\mathbb{E}X\,\mathbb{E}Y$ 只在独立时无条件成立。\\
\textbf{自检}：$X\sim U(0,1)$，算 $\mathbb{E}X$ 与 $\mathrm{Var}X$；若 $X,Y$ 独立同分布，把 $\mathrm{Var}(X+Y)$ 用 $\mathrm{Var}X$ 表示；
再举一个「不相关但不独立」的例子。

\subsubsection*{条件期望}
\textbf{作用}：知道一部分信息之后，对未知量重新给平均。\\
\textbf{解决什么}：最优预测；把「知道的信息」形式化；贝叶斯与滤波的骨架。\\
\textbf{看到就想到}：「已知……求……的期望」「预测」「不确定性变小」「给定条件下的平均」。\\
\textbf{必须会}：离散版 $\mathbb{E}[X\mid B]=\mathbb{E}[X\mathbf{1}_{B}]/P(B)$；全期望公式；$L^{2}$ 中它就是投影。\\
\textbf{最容易混}：$\mathbb{E}[X\mid Y]$ 是随机变量（$Y$ 的函数），不是一个数。\\
\textbf{自检}：掷两枚骰子，已知点数和为 $8$，求第一枚的期望；再写出全期望公式
$\mathbb{E}X=\mathbb{E}\big[\mathbb{E}[X\mid Y]\big]$。

\subsubsection*{特征函数}
\textbf{作用}：一个分布自带的「傅里叶变换」。\\
\textbf{解决什么}：独立和 $\to$ 逐点相乘；证明中心极限定理；判断两个分布是否相同。\\
\textbf{看到就想到}：「独立随机变量相加」「极限分布」「矩」「唯一决定分布」。\\
\textbf{必须会}：$\varphi(t)=\mathbb{E}[e^{itX}]$；独立时相乘；$\varphi^{(k)}(0)=i^{k}\mathbb{E}X^{k}$；正态的 $\varphi$。\\
\textbf{最容易混}：$\varphi$ 是实变量的复值函数；它可以存在而矩不存在（Cauchy 就是）。\\
\textbf{自检}：写 Bernoulli$(p)$ 与 $\mathrm{Exp}(\lambda)$ 的特征函数；再用 $\varphi''(0)$ 取出各自的二阶矩。

\subsubsection*{收敛的几种说法}
\textbf{作用}：区分「越来越近」的几种不同意思。\\
\textbf{解决什么}：说清一个逼近是「每点都近」还是「整体都近」；极限能不能与积分交换顺序。\\
\textbf{看到就想到}：「趋于」「逼近」「一致」「弱」「依分布」。\\
\textbf{必须会}：数列收敛；函数列逐点收敛 vs 一致收敛（差别在「近的快慢是否对所有点统一」）；
依分布收敛 $\mathbb{E}f(X_{n})\to\mathbb{E}f(X)$。\\
\textbf{最容易混}：逐点收敛推不出一致收敛，也推不出积分收敛。\\
\textbf{自检}：举一个在 $[0,1]$ 上逐点收敛但不一致收敛的函数列（$x^{n}$ 就算一个），
并说明为什么这时 $\lim\int\ne\int\lim$。

\subsection{结构：把对称写成可算的规则}
\label{app:grp}

\subsubsection*{群}
\textbf{作用}：把「对称」写成一套能算的规则。\\
\textbf{解决什么}：复合两次仍是同类变换；找出变换下的不变量；把几何问题变成代数问题。\\
\textbf{看到就想到}：「对称」「旋转 / 平移的复合」「一族变换」「不变量」。\\
\textbf{必须会}：四条定义（封闭、结合、单位元、逆元）、矩阵群、子群；同态认得即可。\\
\textbf{最容易混}：群的元素是变换本身，不是被变换的那些点。\\
\textbf{自检}：验证 $\{1,-1\}$（乘法）与 $\{I,R_{\pi}\}$（矩阵乘法）同构，写出对应表；
再说出 $SO(3)$ 的一个子群与它在 $\mathbb{R}^{3}$ 上保持的一个量。
